\section{Proofs for Variational Reduction and Response}
\label{app:reduction-response}

\subsection{Algebraic reduction}

\begin{proof}[Proof of \cref{thm:algebraic-reduction}]
Fix \(z\in\Zspace\).  Partitioning the composite fiber by its intermediate
visible value gives
\[
\begin{aligned}
  ((\rho\circ\pi)_!E)(z)
  &=\inf_{x:\rho(\pi(x))=z}E(x)\\
  &=\inf_{y:\rho(y)=z}\inf_{x:\pi(x)=y}E(x)\\
  &=(\rho_!(\pi_!E))(z).
\end{aligned}
\]
Likewise,
\[
\begin{aligned}
  \inf_x\{E(x)+\ell(\pi(x))\}
  &=\inf_y\inf_{\pi(x)=y}\{E(x)+\ell(y)\}\\
  &=\inf_y\{(\pi_!E)(y)+\ell(y)\}.
\end{aligned}
\]
The base-potential law follows because \(\varphi(\pi(x))=\varphi(y)\) is
constant on each fiber.  The constraint law follows from the definition of
the indicator energy.  No attainment assumption is used.
\end{proof}

\subsection{Existence and uniqueness of the canonical lift}

\begin{proof}[Proof of \cref{thm:canonical-lift}]
Fix \(y\) with finite reduced energy.  Coercivity places every minimizing
sequence in a bounded sublevel set.  Properness makes its closure compact, and
lower semicontinuity gives a minimizer in the closed fiber.

If \(x_0\ne x_1\) were two minimizers, their fiber geodesic would satisfy
\[
  E(x_0\#_{1/2}x_1)
  \le\frac12E(x_0)+\frac12E(x_1)
  -\frac\mu8d(x_0,x_1)^2,
\]
contradicting minimality.  Thus the minimizer is unique.

For equivariance, \(g s_E(y)\) lies in the fiber over \(gy\) and has the same
energy as \(s_E(y)\).  Uniqueness on that fiber gives
\(s_E(gy)=g s_E(y)\).
\end{proof}

\subsection{Smooth minimizer response}

\begin{proof}[Proof of \cref{thm:response-schur}]
We first show continuity of the global minimizer section.  Let \(z_k\to z\).
Choose one fixed trial point near \(\sigma_\star(z)\).  Continuity of \(E\)
uniformly bounds the optimal values for all sufficiently large \(k\).  Local
uniform inf-compactness therefore places \(\sigma_\star(z_k)\) in a common
relatively compact set.  Every convergent subsequence has a limit
\(\bar\sigma\).  Optimality and continuity give, for every trial \(\sigma\),
\[
  E(z,\bar\sigma)
  =\lim E(z_k,\sigma_\star(z_k))
  \le\lim E(z_k,\sigma)=E(z,\sigma).
\]
The minimizer at \(z\) is unique, so \(\bar\sigma=\sigma_\star(z)\).  Every
subsequence has the same limit, proving continuity.

The first-order condition is
\begin{equation}
  E_\sigma(z,\sigma_\star(z))=0.
  \label{eq:appendix-stationarity}
\end{equation}
Its derivative with respect to \(\sigma\) is the invertible matrix \(H\).
The implicit-function theorem gives a local \(C^r\) critical section.  The
continuous global minimizer section agrees with it near every point and is
therefore \(C^r\).  Differentiating
\cref{eq:appendix-stationarity} gives
\[
  E_{\sigma z}+H D_z\sigma_\star=0,
\]
which proves \cref{eq:hidden-response}.  The chain rule gives
\[
  D_z\bar E
  =E_z+E_\sigma D_z\sigma_\star=E_z,
\]
and a second derivative gives
\[
  D^2_{zz}\bar E
  =E_{zz}+E_{z\sigma}D_z\sigma_\star
  =E_{zz}-E_{z\sigma}H^{-1}E_{\sigma z}.
\]

For nested variables \(\sigma=(\alpha,\beta)\), let \(q\) be the quadratic
form defined by the complete Hessian at the optimum.  Positive definiteness of
the joint hidden block makes both quadratic minimizations unique.  Algebraic
reduction gives
\[
  \inf_{\delta\alpha,\delta\beta}
  q(\delta z,\delta\alpha,\delta\beta)
  =\inf_{\delta\alpha}
  \left(\inf_{\delta\beta}
  q(\delta z,\delta\alpha,\delta\beta)\right).
\]
The matrix of a partially minimized quadratic form is the corresponding Schur
complement.  Equality for every \(\delta z\) proves
\cref{eq:schur-associativity}.
\end{proof}

\subsection{Interaction curvature and finite contrasts}

\begin{proof}[Proof of \cref{thm:interaction-curvature}]
For \cref{eq:affine-interventions},
\[
  E_{uu}=0,
  \qquad
  E_{\sigma u}=G.
\]
The \(u\)-block of \cref{eq:schur-effective-hessian} is therefore
\[
  D^2_{uu}\bar E=-G^*H^{-1}G.
\]
For every \(a\in\mathbb R^p\),
\[
  \ip{a}{D^2_{uu}\bar E\,a}
  =-\ip{Ga}{H^{-1}Ga}\le0,
\]
because \(H^{-1}\succ0\).  Expanding in coordinate vectors gives
\cref{eq:pair-interaction-curvature}.  A mixed entry is zero exactly when the
two corresponding covectors are orthogonal under the inner product induced by
\(H^{-1}\).  Finally, without differentiability,
\[
  \bar E(y,u)=\inf_\sigma
  \left(E_0(y,\sigma)+\sum_i u_i\Delta_i(y,\sigma)\right)
\]
is a pointwise infimum of affine functions of \(u\), hence concave.
\end{proof}

\subsection{Finite contrasts and order of reduction}
\label{app:finite-contrasts}

\begin{corollary}[Factorial finite differences]
\label{cor:factorial-contrast}
Assume \([0,1]^p\subset\Uspace\).  For \(T\subseteq[p]\), define
\begin{equation}
  \delta_T\bar E
  =\sum_{B\subseteq T}(-1)^{|T|-|B|}
  \bar E(y,\mathbf1_B).
  \label{eq:factorial-contrast}
\end{equation}
If \(\bar E\in C^{|T|}\), then
\begin{equation}
  \delta_T\bar E
  =\int_{[0,1]^T}
  \partial_T^{|T|}\bar E(u_T,0_{T^c})\,\dd u_T.
  \label{eq:factorial-integral}
\end{equation}
\end{corollary}

\begin{proof}
For a single coordinate, the fundamental theorem of calculus gives
\[
  \bar E(e_i)-\bar E(0)
  =\int_0^1\partial_{u_i}\bar E(t e_i)\,\dd t.
\]
Apply the same identity successively to each coordinate in \(T\).  The product
of difference operators expands as
\[
  \prod_{i\in T}(\mathsf T_i-I)\bar E(0)
  =\sum_{B\subseteq T}(-1)^{|T|-|B|}\bar E(\mathbf1_B),
\]
and repeated integration gives
\[
  \prod_{i\in T}(\mathsf T_i-I)\bar E(0)
  =\int_{[0,1]^T}\partial_T^{|T|}\bar E(u_T,0_{T^c})\,\dd u_T.
\]
For \(T=\{i,j\}\), substitute
\cref{eq:pair-interaction-curvature} to obtain
\cref{eq:pair-factorial-integral}.
\end{proof}

Reduction need not commute with finite differencing.  For
\[
  E_0(\sigma)=\sigma^2,\qquad E_1(\sigma)=(\sigma-1)^2,
\]
both reduced minima are zero, but
\(\inf_\sigma(E_1-E_0)=\inf_\sigma(1-2\sigma)=-\infty\).
Thus hidden states must be reduced first.  If every perturbation
\(\Delta_i\) is constant on each fiber, the minimizer is unchanged and all
higher-order contrasts vanish; this is the exact commuting case.
